\subsection{Operating Regimes of One Model (RQ1)}
We recompute topology--time, topology--geometry, and reference-guided layouts at topology allowances 0, 2\% and 10\% on the same five planar sequences. This produces 3070 layouts from 614 frames. All five settings share each frame's candidate pool, widths, canvas, scalar range, source-value renderer and convex QP implementation. They differ only in the controls in Table~\ref{tab:profiles}; the main regime comparison disables the reference cap ($B=\infty$). A separate Ring grid fixes the reference objective and evaluates all nine combinations of $\varepsilon\in\{0,0.02,0.1\}$ and $B\in\{\infty,0.1,0.05\}$.

The common-model search covers all permutations on 496 frames; the other 118 frames use recorded candidate pools. Across the five settings it evaluates 15141 admissible fixed-order QPs. These coverage counts describe the new QP experiment, not the independently exhaustive frontier experiment below. Selected continuous layouts satisfy their topology and position caps, and the largest normalized candidate-objective gap bound is 1.2e-06. Independent saved-witness checks reconstruct the objective, constraints and selected-order first-order gaps.

Geometry activation reduces the geometry residual relative to the topology--time setting on all five tested inputs. Reference activation greatly reduces absolute reference error, while changing the temporal target from stationary alignment to projected-motion following. These changes express distinct priorities rather than a universal quality ranking. The 2\% and 10\% scalar allowances further reduce pooled reference error on all five sequences in this run, but no monotonic improvement in every component or recomputed temporal history is claimed. Table~\ref{tab:unified_regimes} reports all component measures; Figures~\ref{fig:unified_regimes}--\ref{fig:unified_budget_grid} show the resulting maps.

The shared-input native TMTM and ST-MTM reimplementations are run independently on Ring and Storms. Descriptive similarity uses labelled-pair order agreement and anchor RMSE after one sequence-wide affine gauge fit, including reflection. The fit does not modify the underlying layouts or any reference-error metric. Table~\ref{tab:unified_similarity} reports the results. The topology--time regime is an anchor-alignment surrogate for TMTM, and the topology--geometry regime uses common interval rules instead of ST-MTM's native allocation. In particular, low Storms agreement with native ST-MTM exposes a substantive difference; common design priorities do not imply universally close coordinates. Independent baseline audits retain their separate evidence role in Supplement S7.

The unit tests also isolate an important control: a zero reference weight can still leave a reference-dependent hard position cap. Explicitly disabling that cap recovers geometry-only behavior on a two-feature example, and the canonical null-space rule resolves free newborn coordinates without changing the primary objective. These checks support the interpretation of the parameter regimes rather than native-algorithm equivalence.

\begin{table*}[t]
\caption{New common-model results on all 614 planar frames, with five settings each (3070 layouts). All numbers are percentages: pooled absolute reference error (R-MAE), root mean framewise geometry/temporal squared residual (G-RMS, stationary S-RMS and projected-motion M-RMS), and worst continuous labelled merge error (Top-max). Geometry and temporal quantities use the common fixed canvas, not native-baseline scales. Empty pair/match terms contribute zero. First-frame temporal terms are zero. Different components express competing priorities; lower reference error need not improve every score.}\label{tab:unified_regimes}
\centering\small
\begin{tabular}{llrrrrr}\toprule
Input & Regime & R-MAE & G-RMS & S-RMS & M-RMS & Top-max\\\midrule
Ring & Topology--time & 27.37 & 18.15 & 0.28 & 2.64 & 0.00 \\
Ring & Topology--geometry & 29.75 & 5.88 & 2.66 & 3.67 & 0.00 \\
Ring & Reference, 0\% & 2.57 & 8.58 & 3.97 & 3.42 & 0.00 \\
Ring & Reference, 2\% & 2.28 & 8.41 & 3.33 & 2.66 & 1.49 \\
Ring & Reference, 10\% & 1.17 & 8.15 & 2.65 & 1.20 & 9.40 \\
Storms & Topology--time & 25.31 & 47.41 & 0.38 & 18.65 & 0.00 \\
Storms & Topology--geometry & 33.56 & 11.73 & 8.86 & 19.46 & 0.00 \\
Storms & Reference, 0\% & 6.09 & 20.68 & 20.67 & 9.50 & 0.00 \\
Storms & Reference, 2\% & 5.78 & 20.61 & 20.67 & 8.51 & 1.31 \\
Storms & Reference, 10\% & 5.06 & 19.91 & 20.04 & 6.82 & 8.51 \\
Wildfires & Topology--time & 16.67 & 29.88 & 1.02 & 2.83 & 0.00 \\
Wildfires & Topology--geometry & 13.94 & 8.62 & 7.05 & 7.07 & 0.00 \\
Wildfires & Reference, 0\% & 3.19 & 10.03 & 6.89 & 6.13 & 0.00 \\
Wildfires & Reference, 2\% & 2.72 & 9.65 & 6.26 & 5.29 & 1.85 \\
Wildfires & Reference, 10\% & 1.84 & 9.01 & 5.16 & 4.13 & 8.35 \\
OLR & Topology--time & 24.43 & 35.53 & 0.56 & 9.42 & 0.00 \\
OLR & Topology--geometry & 42.71 & 4.26 & 4.94 & 11.63 & 0.00 \\
OLR & Reference, 0\% & 1.31 & 4.93 & 9.14 & 2.45 & 0.00 \\
OLR & Reference, 2\% & 1.17 & 4.80 & 9.11 & 2.37 & 1.79 \\
OLR & Reference, 10\% & 1.04 & 4.52 & 9.06 & 2.28 & 5.15 \\
PM$_{2.5}$ & Topology--time & 21.71 & 21.60 & 0.74 & 21.88 & 0.00 \\
PM$_{2.5}$ & Topology--geometry & 25.03 & 6.40 & 4.38 & 22.24 & 0.00 \\
PM$_{2.5}$ & Reference, 0\% & 3.05 & 10.60 & 22.07 & 5.44 & 0.00 \\
PM$_{2.5}$ & Reference, 2\% & 2.43 & 9.93 & 21.98 & 4.40 & 1.81 \\
PM$_{2.5}$ & Reference, 10\% & 2.15 & 9.61 & 21.80 & 3.54 & 4.67 \\
\bottomrule\end{tabular}\end{table*}

\begin{table}[t]
\caption{Descriptive similarity of inspired regimes to their native comparison methods. One sequence-wide least-squares affine transform, including reflection, aligns native anchors to the inspired anchors; no framewise fit is used. Agreement counts non-tied labelled feature pairs. RMSE is a percentage of the common inspired canvas. Values measure similarity, not algorithm recovery.}\label{tab:unified_similarity}
\centering\small
\begin{tabular}{llrrr}\toprule
Input & Native & Order (\%) & RMSE (\%) & Pairs\\\midrule
Ring & TMTM & 80.6 & 1.37 & 346 \\
Ring & ST-MTM & 77.2 & 9.69 & 346 \\
Storms & TMTM & 72.5 & 5.98 & 596 \\
Storms & ST-MTM & 51.7 & 34.24 & 596 \\
\bottomrule\end{tabular}\end{table}


\subsection{True Error Budgets and Inverse Queries (RQ2)}\label{sec:budget_experiments}
We newly evaluate nine normalized topology allowances $\{0,0.005,0.01,0.02,0.05,0.1,0.2,0.5,1\}$ on the unchanged five-sequence input population of 614 frames. Every frame uses the same fixed reference, canvas and absolute width convention across the sweep. All permutations are evaluated up to $8!=40320$ orders: 603 frames are exhaustive. The remaining 11 Wildfires frames use the fixed candidate set in Section~\ref{sec:hard_budget}; their values are feasible upper bounds, not certified global minima. No frames are excluded from the reported means.

The run constructs 5526 LP-anchor/barrier witnesses and evaluates 3684 inverse queries at six position caps. Every returned continuous witness meets its declared scalar allowance. No frame's selected position cost increases as $\varepsilon$ increases. The inverse query reports failure when no admissible order meets the requested position cap; on a candidate-only frame this means that no witness was found, rather than global infeasibility. Exhaustive frames supply the stronger interpretation.
\begin{table*}[t]
\caption{New hard-budget experiments on all 614 planar temporal frames. Position columns report the temporal mean of the framewise minimax absolute position error, in percent of the fixed reference extent. These are not NROE or the achieved secondary-QP errors in the retained tables. Wildfires includes 11 candidate-only upper bounds. The inverse column counts frames with a witnessed solution at a 2\% position cap and unrestricted scalar allowance; it is not the count attainable at a 2\% topology cap.}\label{tab:budget_results}
\centering\small
\begin{tabular}{lrrrrr}\toprule
Input & Exact/total & $\varepsilon=0$ & $\varepsilon=2\%$ & $\varepsilon=10\%$ & $E_t(2\%)$ feasible\\\midrule
Storms & 118/118 & 4.326 & 3.499 & 1.889 & 115/118 \\
Wildfires & 19/30 & 6.178 & 3.482 & 1.926 & 30/30 \\
OLR & 60/60 & 0.960 & 0.639 & 0.332 & 60/60 \\
PM$_{2.5}$ & 366/366 & 1.282 & 0.533 & 0.121 & 358/366 \\
Ring & 40/40 & 2.876 & 1.688 & 0.112 & 40/40 \\
\bottomrule\end{tabular}
\end{table*}


At a genuine 2\% topology cap, mean minimax position error changes from 4.326\% to 3.499\% on Storms, 6.178\% to 3.482\% on Wildfires, 0.960\% to 0.639\% on OLR, 1.282\% to 0.533\% on annual PM2.5, and 2.876\% to 1.688\% on Ring. These quantities are continuous per-frame minimax values, not NROE and not the achieved coordinates of the retained secondary-QP layouts. Figure~\ref{fig:budget_response} shows the budget response and illustrative Ring maps generated from the saved witnesses.

At a requested position cap of 2\% with unrestricted scalar allowance, the system finds witnesses in 115/118 Storms, 30/30 Wildfires, 60/60 OLR, 358/366 PM2.5 and 40/40 Ring frames. The three unsuccessful Storms and eight PM2.5 frames were exhaustively searched: their fixed widths and canvas prevent the target even after releasing topology. The inverse query returns a different minimum topology allowance for each feasible frame; these counts do not imply that a 2\% topology cap suffices. A sequence-wide feasibility claim uses the worst frame, whereas Table~\ref{tab:budget_results} reports temporal means.

\paragraph{Scope and controls.}
These experiments establish the continuous non-strict framewise budget behavior. They do not reuse the earlier motion scores as measurements of the new LP witnesses, establish a temporal-QP optimum, or certify a rounded raster under the same cap. The Ring illustrations have finite display resolution and are not used for the numerical guarantees. Existing P/R results below remain separate temporal-construction operating points; their parameter $\kappa$ is an eligibility threshold rather than $\varepsilon$.
